\section{About the Topic}

\subsection{Key Terms}

\textbf{Sovereign debt} is money borrowed by a national government. Its lenders, called creditors, are either multilateral institutions (such as the World Bank and the IMF), bilateral creditors (other governments, for example, China, France, or Japan), or private creditors (commercial banks and bondholders, the investors who buy a government’s bonds). This committee focuses on external debt, meaning what a government owes to foreign creditors.

\textbf{Debt service} is the recurring cost of carrying debt: the interest and principal that fall due each year. It is debt service, not the headline size of the debt, that squeezes budgets year to year, because it must be paid out of the annual budget in direct competition with everything else a government needs to fund.

\textbf{Debt sustainability} describes the ability of a country to meet its current and future debt obligations without requiring unrealistic fiscal adjustment or sacrificing essential public spending.

\textbf{Refinancing / rollover risk.} Governments rarely repay debt out of their cash reserves. When bonds mature, they issue new ones to refinance. Therefore, debt borrowed cheaply a decade ago must be refinanced at today’s interest rates. This means that a bond issued at 4\% can mature into an 8\% market, leading to refinancing at more expensive terms for the government. This “maturity wall” explains why many developing-country debt burdens jumped after 2022.

\textbf{Default vs. debt distress.} A default is the event of missing a payment or formally failing to honour an obligation. Debt distress is the broader, more common condition that the World Bank and IMF grade on a four-band scale, from low risk to in debt distress (International Monetary Fund, 2025). Many countries, although not in outright default, can only honour their obligations by cutting into healthcare, education, and investment spending.

\textbf{Debt restructuring.} In the case of a default, a country renegotiates the terms with its creditors. Negotiations may lead to an extension of the repayment period, a reduction in the interest rate, or a write-down of part of the principal. While restructuring does not erase the debt, it reshapes it into something more manageable for the country. Because they require creditor agreement, restructurings can be complex and take many years.

\textbf{Concessional vs. non-concessional finance.} Concessional finance is lending on terms softer than any market would offer. This can include below-market interest, long repayment periods, and extended grace periods before repayment begins. Grants do not have to be repaid at all.

\textbf{The Paris Club.} An informal group of mostly Western creditor governments that has, since 1956, coordinated the rescheduling or reduction of debts owed to them. It has no treaty or headquarters and runs on the core principle of comparable treatment: a debtor is expected to obtain equally good terms from all its other creditors.

\textbf{The G20 Common Framework.} Agreed in November 2020, this is the main mechanism for restructuring the debts of the poorest countries. It was meant to improve on the Paris Club by bringing in major emerging creditors, including China, India, and Saudi Arabia.

\subsection{Introduction to the Topic}

Although sovereign debt distress is not a new phenomenon, the modern story of this committee is best understood through a clear causal chain. Through the 2010s, money was historically cheap: interest rates in the wealthy parts of the world were near zero, and investors, hunting for higher returns, lent eagerly to developing countries. Governments across Africa, Asia, and Latin America borrowed heavily to fund development. Much of this borrowing increasingly came from a more diverse range of creditors, including bilateral lenders and international capital markets, rather than primarily from traditional concessional sources.\footnote{World Bank. (2022b, December 6). Debt-service payments put biggest squeeze on poor countries since 2000. \url{https://www.worldbank.org/en/news/press-release/2022/12/06/debt-service-payments-put-biggest-squeeze-on-poor-countries-since-2000}} Borrowing continued to rise during the COVID-19 pandemic, as countries sought additional financing to survive the crisis. At the beginning of 2022, central banks raised interest rates sharply to fight inflation that emerged from the pandemic and the energy shock following the Russia-Ukraine War.\footnote{Mosley, L., \& Rosendorff, B. P. (2023). The unfolding sovereign debt crisis. Current History, 122(840), 9–14. \url{https://doi.org/10.1525/curh.2023.122.840.9}} When developing countries’ relatively cheap debt came due and had to be refinanced at the new, higher rates, their debt-service costs rose sharply. What makes this a World Bank problem is that when a country struggles to service its debt or falls into default, its ability to invest in development is directly threatened, which is exactly what the World Bank exists to protect.\footnote{World Bank. (2026, February 17). Beyond the narratives: What the IDR 2025 data reveal about debt in low- and middle-income countries. \url{https://blogs.worldbank.org/en/opendata/beyond-the-narratives--what-the-idr-2025-data-reveal-about-debt-}}

In 2024, developing countries paid a record US\$415 billion in interest alone, resources that could otherwise have gone to schooling, primary healthcare, and essential infrastructure.\footnote{World Bank. (2025b, December 3). Developing countries’ debt outflows hit 50-year high during 2022-2024. \url{https://www.worldbank.org/en/news/press-release/2025/12/03/developing-countries-debt-outflows-hit-50-year-high-during-2022-2024}} Between 2022 and 2024, developing countries paid out US\$741 billion more in debt service than they received in new financing, marking the largest such net outflow in over 50 years.\footnote{World Bank. (2025b, December 3). Developing countries’ debt outflows hit 50-year high during 2022-2024. \url{https://www.worldbank.org/en/news/press-release/2025/12/03/developing-countries-debt-outflows-hit-50-year-high-during-2022-2024}}

In contrast to the IMF, the World Bank is not the institution that runs emergency bailouts. It is the one that finances these countries’ development over decades, and it is now the single-largest provider of net new financing to the 78 most vulnerable countries.\footnote{World Bank. (2025b, December 3). Developing countries’ debt outflows hit 50-year high during 2022-2024. \url{https://www.worldbank.org/en/news/press-release/2025/12/03/developing-countries-debt-outflows-hit-50-year-high-during-2022-2024}} Its relationship to the debt crisis is threefold. First, the Bank is a major creditor and actively varies the terms of its lending based on a country’s debt risk. Second, it is a standard-setter, co-author of the world’s debt-risk scoring system and the leading voice on debt transparency. Third, when distress hits, it helps prevent a debt crisis from becoming a development collapse.\footnote{International Development Association. (2023). Debt. World Bank. \url{https://ida.worldbank.org/en/financing/debt}} \footnote{World Bank. (n.d.-d). Debt (topic overview). \url{https://www.worldbank.org/en/topic/debt/overview}}

\subsection{History of the Topic}

The Bank’s role in sovereign debt crises has evolved through several major episodes. A key starting point is the Latin American debt crisis of the 1980s, when a wave of developing countries (Mexico, Brazil, Argentina, and many others) found themselves unable to service debts accumulated in the 1970s after global interest rates spiked. The decade that followed became known as the “lost decade” for development, because years of austerity and debt repayment left many economies stagnant.\footnote{Federal Reserve History. (2013). Latin American debt crisis of the 1980s. Federal Reserve Board. \url{https://www.federalreservehistory.org/essays/latin-american-debt-crisis}}

In 1996, the World Bank, the IMF, and regional development banks launched the Heavily Indebted Poor Countries (HIPC) Initiative, the first programme in which multilateral institutions agreed to cancel a portion of the debts owed to them by the poorest countries.\footnote{World Bank. (n.d.-b). Heavily Indebted Poor Countries (HIPC) Initiative. \url{https://www.worldbank.org/en/topic/debt/brief/hipc}} HIPC was strengthened in 1999, and in 2005 the G8 went further, agreeing to the Multilateral Debt Relief Initiative (MDRI), which cancelled 100\% of eligible debts owed to IDA, the IMF, and the African Development Fund by countries that completed the HIPC process.\footnote{World Bank. (n.d.-c). Debt relief. \url{https://www.worldbank.org/en/topic/debt-relief}} Together, HIPC and MDRI have relieved 37 participating countries (31 of them in Africa) of more than US\$100 billion in debt, explicitly to free up money for poverty reduction.\footnote{World Bank. (n.d.-b). Heavily Indebted Poor Countries (HIPC) Initiative. \url{https://www.worldbank.org/en/topic/debt/brief/hipc}} This history shows the Bank can act on debt and offers both a model and a warning, as several countries that received relief borrowed their way back into distress within fifteen years.

The relative calm after HIPC paved the way for a new borrowing boom after 2008, when cheap money and new lenders transformed the creditor landscape. Crucially, the share of debt owed to the traditional, coordinated Paris Club creditors fell from 28\% to 11\% by 2020, while debt owed to non-Paris Club bilateral lenders such as China rose from 2\% to 18\%.\footnote{Observer Research Foundation. (2022, June 28). G20’s Debt Service Suspension Initiative: A historical comparison. \url{https://www.orfonline.org/expert-speak/g20s-debt-service-suspension-initiative}} At the same time, borrowing from private bondholders also increased. This shift, from a small club of creditors with a framework to restructure debt together to a fragmented crowd, is one of the reasons today’s crisis is so hard to resolve.

The most recent chapter is the pandemic. In 2020, at the urging of the World Bank and IMF, the G20 launched the Debt Service Suspension Initiative (DSSI), which temporarily suspended debt-service payments for the poorest countries so they could spend on the pandemic response. Between May 2020 and December 2021, the DSSI suspended US\$12.9 billion in payments for 48 of 73 eligible countries.\footnote{World Bank. (2022a). Debt Service Suspension Initiative. \url{https://www.worldbank.org/en/topic/debt/brief/covid-19-debt-service-suspension-initiative}} But it revealed one of the central issues of the new creditor landscape: of all the world’s private creditors, only one agreed to participate, meaning public money freed up by official creditors effectively flowed straight to private bondholders. For more severe cases, the G20 created the Common Framework for Debt Treatments in late 2020, which sought to bring together Paris Club members and the major emerging creditors.

\begin{guidetable}[Y{0.73}Y{0.53}Y{1.74}]{3}
  \textbf{Landmark / Milestone} & \textbf{Year} & \textbf{Summary} \\
  \midrule
  Latin American Debt Crisis & 1980s & A wave of developing country defaults, triggered by spiking global interest rates, led to years of austerity and stagnant growth, known as the “lost decade.” The World Bank and IMF made loans conditional on economic reforms, an approach that drew heavy criticism. \\
  HIPC Initiative launched (World Bank \& IMF) & 1996 & The first programme in which multilateral creditors agreed to cancel portions of their own claims. By 2005, HIPC had provided relief to 37 countries, mostly in Africa, with conditions tied to poverty-reduction strategies. \\
  Multilateral Debt Relief Initiative (MDRI) & 2005 & G8 nations agreed to cancel 100\% of eligible debts owed to IDA, the IMF, and the African Development Fund for HIPC graduates, freeing resources for the UN Millennium Development Goals. \\
  Post-2008 borrowing boom \& new creditors & 2010s & Near-zero global interest rates and eager new lenders, above all China via the Belt and Road Initiative, fuelled a surge in developing-country borrowing. The creditor landscape shifted dramatically away from the Paris Club. \\
  COVID-19 pandemic \& DSSI & 2020–2021 & The G20 launched the DSSI, suspending US\$12.9 billion in debt-service payments for 48 countries. Private creditors almost entirely refused to participate. The G20 Common Framework was created to handle more severe cases. \\
  Post-2022 rate shock \& current wave of distress & 2022–present & Sharp interest rate hikes in advanced economies triggered the refinancing crisis described above. Developing countries paid a record US\$415 billion in interest in 2024 alone. The Common Framework has struggled with slow, contested restructurings. \\
\end{guidetable}
